\documentclass{article}
\usepackage{musixtex}
\usepackage{fontspec}
\usepackage{polyglossia}

\setdefaultlanguage{russian}

\title{Демоверсия ОГЭ по музыке (для гиков)}
\author{Разработчики-Музыканты}
\date{2026}

\begin{document}
\maketitle

\section*{Задание 1. Определение аккорда по коду}
Определите тип аккорда, представленного на нотном стане ниже.

\begin{music}
  % 1 инструмент, 1 нотный стан, скрипичный ключ
  \instrumentnumber{1}
  \setstaffs1{1}
  \setclef1{\treble}

  \startpiece
    % \zh - половинная нота БЕЗ продвижения курсора (Zero-width Half note)
    % Пишем аккорд До-мажор (До-Ми-Соль первой октавы)
    % Буквы c, e, g - это высоты нот в синтаксисе MusiXTeX
    \notes \zh c \zh e \ha g \enotes
    \Endpiece
\end{music}

\section*{Задание 2. Исправление багов в мелодии}
Найдите синтаксическую ошибку в ритмическом рисунке такта.

\begin{music}
  \startpiece
    % \qa - четвертная, \ha - половинная. 
    % Сумма длительностей в такте 4/4 должна сойтись!
    \notes \qa c \qa d \ha e \enotes
    \bar
    \notes \qa f \qa g \qa h \qa i \enotes
    \endpiece
\end{music}

\end{document}
